\documentclass[../DoAn.tex]{subfiles}
\begin{document}

\section{Đặt vấn đề}
\label{section:1.1}
Trò chơi chữ là một hình thức học tập và giải trí có khả năng hỗ trợ người chơi mở rộng vốn từ, rèn luyện khả năng ghi nhớ và liên hệ ngữ nghĩa. Với tiếng Việt, dạng trò chơi này có thêm đặc thù riêng do hệ thống dấu thanh, chữ cái có dấu và nhiều cụm từ ghép có khoảng trắng. Nếu chỉ áp dụng trực tiếp cách xử lý của các trò chơi ô chữ tiếng Anh, hệ thống dễ gặp các vấn đề như xếp chữ sai khi từ có dấu, kiểm tra đáp án không linh hoạt, hoặc khó tổ chức dữ liệu từ vựng theo chủ đề phù hợp với người chơi.

Trong bối cảnh đó, đề tài tập trung xây dựng một trò chơi ô chữ tiếng Việt có thể chơi được trên máy tính, trong đó dữ liệu từ vựng được tổ chức theo chủ đề, mỗi chủ đề có nhiều màn chơi và người chơi nhập đáp án bằng tiếng Việt đầy đủ. Sản phẩm hướng tới việc tạo ra một môi trường chơi đơn giản, rõ ràng, có lưu tiến độ theo tài khoản cục bộ và có cơ chế gợi ý để hỗ trợ người chơi khi gặp từ khó.

Về mặt kỹ thuật, bài toán không chỉ dừng ở việc hiển thị một bảng chữ. Hệ thống cần tạo được bảng từ danh sách từ đầu vào, bảo đảm các từ giao nhau hợp lý, kiểm tra đáp án theo quy tắc tiếng Việt, cập nhật trạng thái khi người chơi đoán đúng và lưu lại kết quả sau mỗi màn. Các yêu cầu này làm cho đề tài có sự kết hợp giữa xử lý chuỗi, thuật toán sắp xếp từ trên lưới, thiết kế giao diện trò chơi và quản lý dữ liệu cục bộ.

Về mặt ứng dụng, sản phẩm có thể được xem như một công cụ giải trí kết hợp học từ vựng. Người chơi không chỉ nhớ mặt chữ mà còn phải đọc gợi ý, liên hệ nghĩa và nhập đúng đáp án. Nếu tiếp tục được mở rộng về dữ liệu và giao diện, trò chơi có thể dùng trong hoạt động tự học, luyện từ theo chủ đề hoặc làm hoạt động nhỏ trong lớp học tiếng Việt.

\section{Mục tiêu và phạm vi đề tài}
\label{section:1.2}
Mục tiêu chính của đồ án là thiết kế và xây dựng một nguyên mẫu trò chơi ô chữ tiếng Việt trên nền tảng Godot. Ứng dụng cần cho phép người chơi chọn tài khoản, chọn chủ đề, chọn màn chơi, xem các gợi ý, nhập đáp án, sử dụng sao để mở gợi ý, nhận điểm khi trả lời đúng và lưu tiến độ hoàn thành theo từng màn. Bên cạnh đó, hệ thống cần có bộ sinh bảng ô chữ tự động từ danh sách từ đã chuẩn hóa để giảm thao tác tạo màn thủ công.

Phạm vi của đồ án tập trung vào chế độ chơi một người trên máy cục bộ. Dữ liệu từ vựng được nhúng trong mã nguồn theo các chủ đề có sẵn; hệ thống chưa hướng tới chế độ nhiều người chơi, đồng bộ dữ liệu qua mạng, cửa hàng nội dung hoặc công cụ biên tập màn chơi dành cho người dùng cuối. Đồ án cũng không đặt mục tiêu xây dựng một bộ sinh ô chữ tối ưu tuyệt đối, mà ưu tiên giải pháp đủ ổn định, dễ hiểu và phù hợp với quy mô của một trò chơi học tập nhỏ.

Các chức năng được ưu tiên trong phạm vi đồ án gồm: quản lý tài khoản cục bộ, chọn chủ đề, chọn màn chơi, sinh bảng ô chữ, hiển thị gợi ý, nhập đáp án, dùng gợi ý bằng sao, tính điểm, tạo nhiệm vụ phụ, hiển thị kết quả và lưu tiến độ. Các chức năng chưa được đưa vào phạm vi gồm: đăng nhập trực tuyến, chia sẻ bảng chơi, thi đấu nhiều người, bảng xếp hạng trực tuyến, công cụ tạo chủ đề bằng giao diện và phát hành chính thức trên kho ứng dụng.

Đề tài cũng giới hạn ở việc đánh giá sản phẩm trên phương diện chức năng và thiết kế kỹ thuật. Phần khảo sát người dùng quy mô lớn, đo lường hiệu quả học tập hoặc so sánh định lượng với các trò chơi khác chưa nằm trong phạm vi hoàn thiện của bản hiện tại. Những nội dung này được xem là hướng phát triển sau khi sản phẩm ổn định hơn.

\section{Định hướng giải pháp}
\label{section:1.3}
Đề tài lựa chọn Godot làm nền tảng phát triển vì Godot hỗ trợ xây dựng giao diện 2D, quản lý scene, tín hiệu và mã GDScript tương đối gọn cho các trò chơi độc lập \cite{godotDocs}. Về xử lý dữ liệu tiếng Việt, hệ thống chuẩn hóa từ nhập vào, tách từ thành các ký tự dùng để xếp bảng, loại bỏ khoảng trắng khi đặt từ trên lưới nhưng vẫn giữ dạng từ đầy đủ để hiển thị và kiểm tra đáp án.

Giải pháp được chia thành các thành phần chính. Bộ cung cấp màn chơi quản lý danh sách chủ đề, tạo các mục từ vựng và chia chúng thành các màn nhỏ. Bộ sinh bảng ô chữ nhận danh sách từ, đặt từ đầu tiên ở gần trung tâm bảng, sau đó lần lượt tìm vị trí giao nhau phù hợp cho các từ còn lại. Giao diện trò chơi hiển thị bảng chữ, danh sách gợi ý, ô nhập đáp án, trạng thái điểm, sao và các hộp thoại hoàn thành. Dữ liệu tiến độ được lưu cục bộ theo từng tài khoản để người chơi có thể quay lại tiếp tục.

Định hướng thiết kế của đồ án là ưu tiên tính rõ ràng và khả năng mở rộng vừa phải. Thay vì xây dựng một thuật toán sinh bảng quá phức tạp, hệ thống sử dụng phương pháp heuristic có thể giải thích được và dễ điều chỉnh. Thay vì nhập từng ký tự vào từng ô, hệ thống cho phép nhập toàn bộ đáp án bằng ô nhập văn bản để tận dụng bộ gõ tiếng Việt quen thuộc của hệ điều hành. Thay vì lưu dữ liệu trong cơ sở dữ liệu lớn, hệ thống dùng JSON cục bộ vì lượng dữ liệu tiến độ nhỏ và không cần máy chủ.

\section{Bố cục đồ án}
\label{section:1.4}
Phần còn lại của báo cáo được tổ chức như sau. Chương 2 trình bày khảo sát hiện trạng, phân tích yêu cầu chức năng và yêu cầu phi chức năng của trò chơi ô chữ tiếng Việt. Chương 3 giới thiệu nền tảng lý thuyết và các công nghệ sử dụng, bao gồm Godot, GDScript, xử lý chuỗi tiếng Việt và mô hình lưu dữ liệu cục bộ. Chương 4 trình bày thiết kế kiến trúc, thiết kế giao diện, thiết kế lớp, xây dựng ứng dụng, kiểm thử và triển khai thử nghiệm. Chương 5 tập trung vào các giải pháp và đóng góp nổi bật của đồ án, đặc biệt là xử lý từ tiếng Việt, sinh bảng ô chữ và tổ chức tiến trình chơi. Chương 6 tổng kết kết quả đạt được, hạn chế còn tồn tại và các hướng phát triển tiếp theo.

\end{document}
